\documentclass[11pt,a4paper]{article}
\usepackage[left=21mm,right=21mm,top=18mm,bottom=20mm]{geometry}
\usepackage[T1]{fontenc}
\usepackage[utf8]{inputenc}
\usepackage{lmodern,amsmath,amssymb,mathtools}
\usepackage{graphicx,booktabs,longtable,array,calc}
\usepackage{microtype,enumitem,needspace,etoolbox,titlesec}
\usepackage{xcolor,xurl,hyperref}
\definecolor{noteblue}{HTML}{1458A6}
\hypersetup{colorlinks=true,urlcolor=noteblue,linkcolor=black,pdftitle={Mathematical
Analysis - Week 1},pdfauthor={Jiayu Chen}}
\setlength{\parindent}{0pt}
\setlength{\parskip}{4pt plus 1pt minus 1pt}
\linespread{1.035}
\setlength{\emergencystretch}{3em}
\setlength{\abovedisplayskip}{7pt plus 2pt minus 1pt}
\setlength{\belowdisplayskip}{7pt plus 2pt minus 1pt}
\setlength{\abovedisplayshortskip}{4pt plus 1pt}
\setlength{\belowdisplayshortskip}{5pt plus 1pt}
\setlist{itemsep=3pt,topsep=4pt,parsep=1pt,leftmargin=1.7em}
\allowdisplaybreaks[1]
\raggedbottom
\clubpenalty=4000
\widowpenalty=4000
\providecommand{\tightlist}{}
\newcommand{\passthrough}[1]{#1}
\newcommand{\pandocbounded}[1]{#1}
\setcounter{secnumdepth}{0}
\setcounter{tocdepth}{1}
\titleformat{\section}{\Large\bfseries}{\thesection}{0pt}{}
\titleformat{\subsection}{\large\bfseries}{\thesubsection}{0pt}{}
\titleformat{\subsubsection}{\normalsize\bfseries}{\thesubsubsection}{0pt}{}
\titlespacing*{\section}{0pt}{14pt}{8pt}
\titlespacing*{\subsection}{0pt}{12pt}{6pt}
\titlespacing*{\subsubsection}{0pt}{9pt}{4pt}
\pretocmd{\section}{\Needspace{6\baselineskip}}{}{}
\pretocmd{\subsection}{\Needspace{5\baselineskip}}{}{}
\pretocmd{\subsubsection}{\Needspace{4\baselineskip}}{}{}
\makeatletter
\renewcommand*\l@section{\@dottedtocline{1}{0em}{0em}}
\makeatother
\pagestyle{plain}
\begin{document}
\begin{center}
{\LARGE Mathematical Analysis - Week 1\par}
\vspace{5pt}
{\large Lectures 1-4\par}
\end{center}
{\small\tableofcontents}
\vspace{7pt}
The first question is not how to substitute a number into a formula, but
what the values of the formula do near that number. A function can have
a limit at a point even when its value there is missing or different. We
start with that distinction, make it precise, and then use it to study
continuity.

\section{Preliminary}\label{preliminary}

\begin{itemize}
\item
  \(\mathbb{N}\) denotes the set of all positive integers.
\item
  \(\mathbb{Z}\) denotes the set of all integers.
\item
  \(\mathbb{Q}\) denotes the set of all rational numbers.
\item
  \(\mathbb{R}\) denotes the set of all real numbers.
\end{itemize}

\section{Lecture 1: approaching a
point}\label{lecture-1-approaching-a-point}

\subsection{1. Key idea}\label{key-idea}

\begin{quote}
We use limits to describe the tendency of a function as the variable
approaches a given point, but not the value at the point.
\end{quote}

\subsection{2. Important examples}\label{important-examples}

\subsubsection{Example 1: Same limit, different value at the
point}\label{example-1-same-limit-different-value-at-the-point}

\[
f(x)=
\begin{cases}
x^2\sin\dfrac{1}{x},&x\ne0,\\
0,&x=0,
\end{cases}
\]

\includegraphics[width=0.52\linewidth,height=\textheight,keepaspectratio]{image.png}

\[
g(x)=
\begin{cases}
x\sin\dfrac{1}{x},&x\ne0,\\
1,&x=0.
\end{cases}
\]

\includegraphics[width=0.52\linewidth,height=\textheight,keepaspectratio]{image-1.png}

For \(x\ne0\), the useful estimates are

\[
|f(x)|\le |x|^2,\qquad |g(x)|\le |x|.
\]

Therefore:

\[
\lim_{x\to0}f(x)=0,\qquad
\lim_{x\to0}g(x)=0.
\]

My takeaway:

\begin{quote}
If \(x\) gets close to 0, keeping \(x\ne0\), \(f(x)\) approaches 0 with
error \(\leq |x|^2\); \(g(x)\) approaches 0 with error \(\leq |x|\).
\end{quote}

\subsubsection{Example 2: average velocity and instantaneous
velocity}\label{example-2-average-velocity-and-instantaneous-velocity}

\textbf{Average velocity:}

\[
\frac{\Delta s}{\Delta t}
=\frac{f(t_2)-f(t_1)}{t_2-t_1}.
\]

With \(h=t_2-t_1\):

\[
\frac{f(t_1+h)-f(t_1)}{h}.
\]

\textbf{Instantaneous velocity at \(t_0\):}

\[
v(t_0)=\lim_{h\to0}\frac{f(t_0+h)-f(t_0)}{h}
=\lim_{t\to t_0}\frac{f(t)-f(t_0)}{t-t_0}.
\]

\subsection{Informal definition of
limit}\label{informal-definition-of-limit}

\begin{quote}
Assume \(f(x)\) is defined on an open interval about \(c\), except
possibly at \(c\) itself. If \(f(x)\) is arbitrarily close to \(L\) for
all \(x\) sufficiently close to \(c\) other than \(c\) itself, then
\end{quote}

\[
\lim_{x\to c}f(x)=L.
\]

This is read ``the limit of \(f(x)\) as \(x\) approaches \(c\) is
\(L\).''

\subsection{3. Limit calculation
techniques}\label{limit-calculation-techniques}

\subsubsection{Algebraic simplification}\label{algebraic-simplification}

\[
\lim_{x\to1}\frac{x^2-1}{x-1}
=\lim_{x\to1}(x+1)
=2.
\]

\subsubsection{Rationalizing the
denominator}\label{rationalizing-the-denominator}

\[
\lim_{x\to1}\frac{\sqrt{x}-1}{x-1}
=\lim_{x\to1}\frac{1}{\sqrt{x}+1}
=\frac12.
\]

\textbf{Remark:} Simplification when calculating limits.

\textbf{General strategy:}

\begin{enumerate}
\def\labelenumi{\arabic{enumi}.}
\item
  Check direct substitution.
\item
  Factor or rationalize if necessary.
\item
  Simplify, then evaluate the limit.
\end{enumerate}

\subsection{\texorpdfstring{Class 2: Describing the behavior of a
function as
\(x\to0\)}{Class 2: Describing the behavior of a function as x\textbackslash to0}}\label{class-2-describing-the-behavior-of-a-function-as-xto0}

\textbf{1. Jump}

\includegraphics[width=0.52\linewidth,height=\textheight,keepaspectratio]{image-2.png}

\textbf{2. Unbounded}

\includegraphics[width=0.52\linewidth,height=\textheight,keepaspectratio]{image-3.png}

\textbf{3. Oscillation}

\includegraphics[width=0.52\linewidth,height=\textheight,keepaspectratio]{image-4.png}

\subsection{4. When a limit does not
exist}\label{when-a-limit-does-not-exist}

\begin{longtable}[]{@{}lll@{}}
\toprule\noalign{}
Type & Behavior near the point & Example \\
\midrule\noalign{}
\endhead
\bottomrule\noalign{}
\endlastfoot
Jump & Different left and right limits & Unit step function \\
Unbounded & Values grow without bound & \(1/x\) near zero \\
Oscillation & No single limiting value & \(\sin(1/x)\) near zero \\
\end{longtable}

\section{Lecture 2: the definition and how to use
it}\label{lecture-2-the-definition-and-how-to-use-it}

\subsection{Definition of limit
(Informal)}\label{definition-of-limit-informal}

\[
g(x)=
\begin{cases}
2x, & \text{if } x\ne4,\\
0, & \text{if } x=4.
\end{cases}
\]

By the informal description of a limit: if \(x\) gets sufficiently close
to \(4\), but is not equal to \(4\), then \(g(x)\) gets arbitrarily
close to \(8\).

We need to show that the gap between \(g(x)\) and \(L=8\), namely
\(|g(x)-8|\), can be made as small as we choose if \(x\) is kept close
enough to \(c=4\), but is not equal to \(4\).

\subsection{Definition}\label{definition}

Let \(f(x)\) be defined on an open interval \(I\) about \(c\), except
possibly at \(c\) itself. We say that the limit of \(f(x)\) as \(x\)
approaches \(c\) is the number \(L\), and write

\[
\lim_{x\to c}f(x)=L,
\]

if, for every \(\varepsilon>0\), there exists a corresponding
\(\delta>0\) such that

\[
|f(x)-L|<\varepsilon
\quad\text{whenever }x\in I\text{ and }0<|x-c|<\delta.
\]

\subsection{Simpler writing}\label{simpler-writing}

We say \(\lim_{x\to c}f(x)=L\) if, for any \(\varepsilon>0\), there
exists \(\delta>0\) such that

\[
\text{if }x\in I,\quad 0<|x-c|<\delta,
\quad\text{then }|f(x)-L|<\varepsilon.
\]

\subsection{Personal interpretation}\label{personal-interpretation}

\begin{itemize}
\item
  \(\varepsilon\): how close we require \(f(x)\) to be to \(L\).
\item
  \(\delta\): how close \(x\) needs to be to \(c\) in order to satisfy
  that requirement.
\item
  \(|f(x)-L|<\varepsilon\): \(f(x)\) lies within an
  \(\varepsilon\)-neighborhood of \(L\).
\item
  \(0<|x-c|<\delta\): \(x\) lies within a \(\delta\)-neighborhood of
  \(c\), but \(x\ne c\).
\end{itemize}

The key idea is

\[
\boxed{\text{Given any }\varepsilon>0,\text{ we can find a }\delta>0}
\]

such that

\[
x\text{ is sufficiently close to }c
\quad\Longrightarrow\quad
f(x)\text{ is within }\varepsilon\text{ of }L.
\]

In other words, \(\varepsilon\) is the required accuracy for the output,
while \(\delta\) is how much accuracy we need for the input.

Thus, the \(\varepsilon\)-\(\delta\) definition gives a precise
mathematical meaning to the intuitive statement: if \(x\) is
sufficiently close to \(c\), then \(f(x)\) is sufficiently close to
\(L\).

\subsection{Useful remarks}\label{useful-remarks}

The order of the quantifiers matters. The number \(\delta\) may depend
on \(\varepsilon\), on the function, and on the point \(c\), but not on
the particular \(x\) that will later be chosen. We need one radius that
works for every point in the punctured interval. Making \(\delta\)
smaller never spoils a valid choice.

The condition \(0<|x-c|\) deliberately excludes \(c\). Continuity will
put the value at \(c\) back into the discussion.

We shall often use the following arguments:

\begin{enumerate}
\def\labelenumi{\arabic{enumi}.}
\item
  Let \(a\) be a real number. If, for every \(\varepsilon>0\), we have
  \(a\leq\varepsilon\), then \(a\leq0\).
\item
  Let \(a\) be a real number. If, for every \(\varepsilon>0\), we have
\end{enumerate}

\[
-\varepsilon\leq a\leq\varepsilon,
\]

then \(a=0\).

\subsection{\texorpdfstring{Example: revisiting
\(g(x)\)}{Example: revisiting g(x)}}\label{example-revisiting-gx}

Consider

\[
g(x)=
\begin{cases}
2x, & \text{if }x\ne4,\\
0, & \text{if }x=4.
\end{cases}
\]

Show that

\[
\lim_{x\to4}g(x)=8.
\]

\textbf{Analysis.} For every \(\varepsilon>0\), we need

\[
|g(x)-8|=|2x-8|<\varepsilon,
\qquad x\ne4,
\]

which requires

\[
0<|x-4|<\frac{\varepsilon}{2}.
\]

\textbf{Proof.} Given arbitrary \(\varepsilon>0\), take

\[
\delta=\frac{\varepsilon}{2}.
\]

Then

\[
|g(x)-8|=|2x-8|=2|x-4|<2\delta=\varepsilon
\]

whenever \(0<|x-4|<\delta\). By the definition of limit,

\[
\lim_{x\to4}g(x)=8.
\]

\subsection{Exercise}\label{exercise}

Prove that

\[
\lim_{x\to1}\frac{1}{x}=1.
\]

\textbf{Hint.} Restrict \(x\) by taking \(x>\frac12\). Then

\[
\left|\frac{1}{x}-1\right|=\frac{|x-1|}{|x|}<2|x-1|.
\]

\textbf{Solution.} Given \(\varepsilon>0\), choose
\(\delta=\min\{1/2,\varepsilon/2\}\). If \(0<|x-1|<\delta\), then
\(x>1/2\) and the displayed estimate is less than \(\varepsilon\). The
first restriction keeps the denominator away from zero; the second
controls the error.

\subsection{Exercise 1.1.5 revisited}\label{exercise-1.1.5-revisited}

Using the definition of limit, prove that the following functions have
no limit as \(x\to0\).

\subsubsection{(a)}\label{a}

\[
U(x)=
\begin{cases}
0, & x<0,\\
1, & x\geq0.
\end{cases}
\]

\subsubsection{(b)}\label{b}

\[
g(x)=
\begin{cases}
\dfrac{1}{x}, & x\ne0,\\
0, & x=0.
\end{cases}
\]

\subsubsection{(c)}\label{c}

\[
f(x)=
\begin{cases}
0, & x\leq0,\\
\sin\dfrac{1}{x}, & x>0.
\end{cases}
\]

\textbf{Tip.} Suppose the limit exists and derive a contradiction.

\subsubsection{For (b)}\label{for-b}

Suppose

\[
\lim_{x\to0}\frac{1}{x}=L.
\]

Use the definition with \(\varepsilon_0=1\). It would give a radius
\(\delta>0\) on which \(|1/x-L|<1\). Choose

\[
0<x<\min\{\delta,1/(|L|+2)\}.
\]

Then \(1/x>|L|+2\), so \(|1/x-L|\ge 1/x-|L|>2\), a contradiction.

\textbf{Solutions to (a) and (c).} For the step function, points on the
two sides have values \(0\) and \(1\). If a limit \(L\) existed, choose
\(\varepsilon=1/3\) and a common radius. These two values would satisfy
\(1\le |1-L|+|L|<2/3\), which is impossible.

For (c), positive points \(x_n=(\pi/2+2n\pi)^{-1}\) and
\(y_n=(3\pi/2+2n\pi)^{-1}\) approach zero and give values \(1\) and
\(-1\). A proposed limit would have to be within \(1/2\) of both values
for all sufficiently large \(n\), again impossible. These sequences are
a convenient way to choose points in every sufficiently small
neighborhood; a formal sequence criterion comes in Week 2.

\section{Lecture 3: limit laws and one-sided
limits}\label{lecture-3-limit-laws-and-one-sided-limits}

\textbf{Date:} 2026.9.10

\subsection{Theorem 1.1.3 (Limit Laws)}\label{theorem-1.1.3-limit-laws}

Assume that

\[
\lim_{x\to c}f(x)=L
\qquad\text{and}\qquad
\lim_{x\to c}g(x)=M.
\]

\subsubsection{(a) Sum Rule}\label{a-sum-rule}

\[
\lim_{x\to c}\bigl(f(x)+g(x)\bigr)
=\lim_{x\to c}f(x)+\lim_{x\to c}g(x).
\]

\textbf{Proof.} Let \(\varepsilon>0\). Choose \(\delta_1,\delta_2>0\) so
that \(|f(x)-L|<\varepsilon/2\) and \(|g(x)-M|<\varepsilon/2\). For
\(\delta=\min\{\delta_1,\delta_2\}\),

\[
|(f(x)+g(x))-(L+M)|\le |f(x)-L|+|g(x)-M|<\varepsilon.
\]

\subsubsection{(b) Product Rule}\label{b-product-rule}

\[
\lim_{x\to c}\bigl(f(x)\cdot g(x)\bigr)
=\lim_{x\to c}f(x)\cdot\lim_{x\to c}g(x).
\]

\textbf{Proof.} Since \(f(x)\to L\), it is bounded near \(c\); take a
neighborhood where \(|f(x)|\le |L|+1\). Then

\[
|f(x)g(x)-LM|\le |f(x)||g(x)-M|+|M||f(x)-L|,
\]

For an explicit choice, first restrict to \(|f(x)-L|<1\). Also impose

\[
|g(x)-M|<\frac{\varepsilon}{2(|L|+1)},
\qquad |f(x)-L|<\frac{\varepsilon}{2(|M|+1)}.
\]

Take the minimum of the three radii. The two terms in the estimate are
then respectively less than \(\varepsilon/2\) and at most
\(\varepsilon/2\). This also covers \(L=0\) or \(M=0\), without division
by either number.

\subsubsection{(c) Quotient Rule}\label{c-quotient-rule}

If \(M\ne0\), then

\[
\lim_{x\to c}\frac{f(x)}{g(x)}
=\frac{\displaystyle\lim_{x\to c}f(x)}
{\displaystyle\lim_{x\to c}g(x)}.
\]

\textbf{Proof.} Since \(g(x)\to M\ne0\), eventually \(|g(x)|>|M|/2\).
Hence

\[
\left|\frac1{g(x)}-\frac1M\right|
=\frac{|g(x)-M|}{|g(x)||M|}
\le \frac{2|g(x)-M|}{|M|^2}\to0.
\]

Thus \(1/g(x)\to1/M\), and the quotient rule follows from the product
rule.

\subsection{Remark: Does the Converse
Hold?}\label{remark-does-the-converse-hold}

The Sum Rule requires that both limits on the right-hand side exist. In
general, the converse is false: the existence of

\[
\lim_{x\to c}(f(x)+g(x))
\]

does not imply that both \(\lim_{x\to c}f(x)\) and \(\lim_{x\to c}g(x)\)
exist.

\subsubsection{Counterexample}\label{counterexample}

Let

\[
f(x)=\sin\frac{1}{x},
\qquad
g(x)=-\sin\frac{1}{x}.
\]

Consider \(x\to0\). Although

\[
\lim_{x\to0}\sin\frac{1}{x}
\]

does not exist, we have

\[
f(x)+g(x)=0.
\]

Therefore,

\[
\lim_{x\to0}(f(x)+g(x))=0
\]

exists. Thus,

\[
\boxed{
\lim(f+g)\text{ exists}
\not\Longrightarrow
\lim f,\lim g\text{ both exist}
}
\]

The problematic behavior of two functions can cancel each other out.

\subsection{Exercise 1.1.9 (Limits of
Polynomials)}\label{exercise-1.1.9-limits-of-polynomials}

If

\[
P(x)=\sum_{k=0}^{n}a_kx^k
=a_nx^n+a_{n-1}x^{n-1}+\cdots+a_0,
\]

then

\[
\lim_{x\to c}P(x)
=\sum_{k=0}^{n}a_kc^k
=P(c).
\]

\subsection{Exercise 1.1.10 (Limits of Rational
Functions)}\label{exercise-1.1.10-limits-of-rational-functions}

If \(P(x)\) and \(Q(x)\) are polynomials and \(Q(c)\ne0\), then

\[
\lim_{x\to c}\frac{P(x)}{Q(x)}=\frac{P(c)}{Q(c)}.
\]

\subsection{Theorem 1.1.4 (The Sandwich
Theorem)}\label{theorem-1.1.4-the-sandwich-theorem}

Suppose that

\[
g(x)\leq f(x)\leq h(x)
\]

for all \(x\) in some open interval containing \(c\), except possibly at
\(x=c\). Suppose also that

\[
\lim_{x\to c}g(x)=\lim_{x\to c}h(x)=L.
\]

Then

\[
\lim_{x\to c}f(x)=L.
\]

\textbf{Proof.} Given \(\varepsilon>0\), choose \(\delta>0\) so that
\(L-\varepsilon<g(x)\) and \(h(x)<L+\varepsilon\) whenever
\(0<|x-c|<\delta\). The inequalities \(g(x)\le f(x)\le h(x)\) then imply
\(|f(x)-L|<\varepsilon\).

\includegraphics[width=0.52\linewidth,height=\textheight,keepaspectratio]{image-5.png}

\subsection{Exercise 1.1.12}\label{exercise-1.1.12}

Assume \(I\) is an open interval, \(c\in I\), and

\[
f(x)\leq b
\qquad\text{for all }x\in I\setminus\{c\}.
\]

Then

\[
\lim_{x\to c}f(x)\leq b,
\]

provided the limit exists.

\textbf{Proof.} Write \(L=\lim_{x\to c}f(x)\). If \(L>b\), take
\(\varepsilon=(L-b)/2\). Near \(c\) the definition gives
\(f(x)>L-\varepsilon=(L+b)/2>b\), contradicting the assumed upper bound.

\subsection{Exercise 1.1.13 (Order-preserving Property of
Limits)}\label{exercise-1.1.13-order-preserving-property-of-limits}

If

\[
f(x)\leq g(x)
\]

for all \(x\) in some open interval containing \(c\), except possibly at
\(x=c\), and both limits exist, then

\[
\lim_{x\to c}f(x)\leq\lim_{x\to c}g(x).
\]

Limits preserve the non-strict inequality \(\leq\), but may not preserve
the strict inequality \(<\). For example, let

\[
f(x)=x^2,
\qquad g(x)=2x^2.
\]

For \(x\ne0\), \(f(x)<g(x)\), but

\[
\lim_{x\to0}x^2=\lim_{x\to0}2x^2=0.
\]

\subsection{Lemma 1.1.5}\label{lemma-1.1.5}

For \(\theta\) measured in radians,

\[
-|\theta|\leq\sin\theta\leq|\theta|,
\qquad
-|\theta|\leq1-\cos\theta\leq|\theta|.
\]

In particular, for \(0<\theta\leq\dfrac{\pi}{2}\),

\[
0<\sin\theta<\theta.
\]

\textbf{Geometric proof.} For \(0<\theta<\pi/2\), take \(A=(1,0)\) and
\(P=(\cos\theta,\sin\theta)\) on the unit circle. The chord \(AP\) is
shorter than the arc of length \(\theta\). Therefore

\[
\sin^2\theta+(1-\cos\theta)^2=|AP|^2<\theta^2.
\]

Each nonnegative term gives \(|\sin\theta|<\theta\) and
\(1-\cos\theta<\theta\) in this range. At \(\theta=\pi/2\) the claimed
inequalities also hold. Symmetry handles negative angles. For
\(|\theta|\ge\pi/2\), use \(|\sin\theta|\le1\le|\theta|\) and
\(1-\cos\theta=2\sin^2(\theta/2)\le 2|\sin(\theta/2)|\le |\theta|\); the
last sine bound is already established globally.

This argument precedes differentiation. Using the mean value theorem
here would reverse the logical order: the derivatives of sine and cosine
will themselves use the trigonometric limit below.

\includegraphics[width=0.52\linewidth,height=\textheight,keepaspectratio]{image-7.png}

\subsection{One-sided Limits}\label{one-sided-limits}

\subsubsection{Left-hand limit}\label{left-hand-limit}

\[
\lim_{x\to c^-}f(x)=L
\]

if, for every \(\varepsilon>0\), there exists \(\delta>0\) such that

\[
c-\delta<x<c
\quad\Longrightarrow\quad
|f(x)-L|<\varepsilon.
\]

\subsubsection{Right-hand limit}\label{right-hand-limit}

\[
\lim_{x\to c^+}f(x)=L
\]

if, for every \(\varepsilon>0\), there exists \(\delta>0\) such that

\[
c<x<c+\delta
\quad\Longrightarrow\quad
|f(x)-L|<\varepsilon.
\]

The two-sided limit exists if and only if the one-sided limits both
exist and are equal:

\[
\lim_{x\to c}f(x)=L
\quad\Longleftrightarrow\quad
\lim_{x\to c^-}f(x)=\lim_{x\to c^+}f(x)=L.
\]

\subsubsection{\texorpdfstring{Examples at
\(x\to0\)}{Examples at x\textbackslash to0}}\label{examples-at-xto0}

\paragraph{\texorpdfstring{Example 1: The \(0\)-\(1\) step
function.}{Example 1: The 0-1 step function.}}\label{example-1-the-0-1-step-function.}

\[
U(x)=
\begin{cases}
0, & x<0,\\
1, & x\geq0.
\end{cases}
\]

Then

\[
\lim_{x\to0^-}U(x)=0,
\qquad
\lim_{x\to0^+}U(x)=1,
\]

so the two-sided limit does not exist.

\paragraph{\texorpdfstring{Example 2: The function
\(1/x\).}{Example 2: The function 1/x.}}\label{example-2-the-function-1x.}

\[
\lim_{x\to0^-}\frac{1}{x}=-\infty,
\qquad
\lim_{x\to0^+}\frac{1}{x}=+\infty.
\]

Thus, \(1/x\) has no finite two-sided limit at \(0\).

\paragraph{\texorpdfstring{Example 3: The oscillating function
\(\sin(1/x)\).}{Example 3: The oscillating function \textbackslash sin(1/x).}}\label{example-3-the-oscillating-function-sin1x.}

\[
\lim_{x\to0^-}\sin\frac{1}{x}
\quad\text{and}\quad
\lim_{x\to0^+}\sin\frac{1}{x}
\]

both do not exist.

\paragraph{\texorpdfstring{The density of \(\mathbb{Q}\) in
\(\mathbb{R}\).}{The density of \textbackslash mathbb\{Q\} in \textbackslash mathbb\{R\}.}}\label{the-density-of-mathbbq-in-mathbbr.}

We say that \(\mathbb{Q}\) is dense in \(\mathbb{R}\) if every non-empty
open interval contains a rational number. For any \(a,b\in\mathbb{R}\)
with \(a<b\), there exists \(q\in\mathbb{Q}\) such that

\[
a<q<b.
\]

Equivalently, for every \(x\in\mathbb{R}\) and every \(\varepsilon>0\),
there exists \(q\in\mathbb{Q}\) such that

\[
|q-x|<\varepsilon.
\]

The irrational numbers are also dense in \(\mathbb{R}\).

\paragraph{Example 4: The Dirichlet
function.}\label{example-4-the-dirichlet-function.}

\[
D(x)=
\begin{cases}
1, & x\in\mathbb{Q},\\
0, & x\in\mathbb{R}\setminus\mathbb{Q}.
\end{cases}
\]

Every interval to the left or right of \(0\) contains both rational and
irrational numbers. Hence,

\[
\lim_{x\to0^-}D(x)
\quad\text{and}\quad
\lim_{x\to0^+}D(x)
\]

both do not exist.

\subsection{Theorem 1.1.6}\label{theorem-1.1.6}

Suppose that \(f(x)\) is defined on an open interval containing \(c\),
except possibly at \(c\) itself. Then \(f(x)\) has a limit as \(x\)
approaches \(c\) if and only if its left-hand and right-hand limits both
exist and are equal:

\[
\lim_{x\to c}f(x)=L
\quad\Longleftrightarrow\quad
\lim_{x\to c^-}f(x)=\lim_{x\to c^+}f(x)=L.
\]

\subsubsection{Proof}\label{proof}

If \(\lim_{x\to c}f(x)=L\), the \(\varepsilon\)-\(\delta\) condition
immediately implies the same estimate for \(x<c\) and for \(x>c\).
Therefore both one-sided limits equal \(L\).

Conversely, suppose both one-sided limits equal \(L\). Given
\(\varepsilon>0\), choose \(\delta_1,\delta_2>0\) from the left- and
right-hand definitions and let

\[
\delta=\min\{\delta_1,\delta_2\}.
\]

If \(0<|x-c|<\delta\), then either \(c-\delta<x<c\) or \(c<x<c+\delta\).
In either case,

\[
|f(x)-L|<\varepsilon.
\]

Thus \(\lim_{x\to c}f(x)=L\).

\subsection{Theorem 1.1.7}\label{theorem-1.1.7}

For \(\theta\) measured in radians,

\[
\lim_{\theta\to0}\frac{\sin\theta}{\theta}=1.
\]

\includegraphics[width=0.52\linewidth,height=\textheight,keepaspectratio]{image-8.png}

\subsubsection{Proof}\label{proof-1}

For \(0<\theta<\dfrac{\pi}{2}\),

\[
\cos\theta\leq\frac{\sin\theta}{\theta}\leq1.
\]

Since \(\lim_{\theta\to0}\cos\theta=1\), the Sandwich Theorem gives

\[
\lim_{\theta\to0^+}\frac{\sin\theta}{\theta}=1.
\]

For \(\theta<0\),

\[
\frac{\sin\theta}{\theta}
=\frac{\sin(-\theta)}{-\theta},
\]

so the left-hand limit is also \(1\). Therefore,

\[
\boxed{\displaystyle\lim_{\theta\to0}\frac{\sin\theta}{\theta}=1}.
\]

\subsubsection{Consequences}\label{consequences}

For any constant \(k\ne0\),

\[
\lim_{\theta\to0}\frac{\sin(k\theta)}{\theta}
=k\lim_{\theta\to0}\frac{\sin(k\theta)}{k\theta}=k.
\]

Also,

\[
\lim_{x\to0}\frac{\sin(\sin x)}{\sin x}=1,
\]

and, for \(n\) nested sine functions,

\[
\lim_{x\to0}
\frac{\sin\bigl(\sin(\cdots\sin x\cdots)\bigr)}{x}=1.
\]

\subsubsection{\texorpdfstring{Proof of
\(\displaystyle\lim_{x\to0}\frac{\sin(\sin x)}{x}=1\)}{Proof of \textbackslash displaystyle\textbackslash lim\_\{x\textbackslash to0\}\textbackslash frac\{\textbackslash sin(\textbackslash sin x)\}\{x\}=1}}\label{proof-of-displaystylelim_xto0fracsinsin-xx1}

Since \(\lim_{x\to0}\sin x=0\), Theorem 1.1.7 gives

\[
\lim_{x\to0}\frac{\sin(\sin x)}{\sin x}=1.
\]

Moreover,

\[
\frac{\sin(\sin x)}{x}
=\frac{\sin(\sin x)}{\sin x}\cdot\frac{\sin x}{x}.
\]

Therefore,

\[
\begin{aligned}
\lim_{x\to0}\frac{\sin(\sin x)}{x}
&=\left(\lim_{x\to0}\frac{\sin(\sin x)}{\sin x}\right)
\left(\lim_{x\to0}\frac{\sin x}{x}\right)\\
&=1\cdot1=1.
\end{aligned}
\]

\section{Lecture 4: continuity}\label{lecture-4-continuity}

\subsection{Example 1.1.19}\label{example-1.1.19}

\[
f(x)=
\begin{cases}
x, & x<0,\\
1, & 0\leq x<1,\\
2, & x=1,\\
x, & x>1.
\end{cases}
\]

\includegraphics[width=0.52\linewidth,height=\textheight,keepaspectratio]{image-9.png}

\subsection{Definition 1.1.8 (Continuity at a
Point)}\label{definition-1.1.8-continuity-at-a-point}

Let \(c\) be a real number that is either an interior point or an
endpoint of an interval in the domain of \(f\).

\subsubsection{(a) Interior point}\label{a-interior-point}

\(f\) is continuous at \(c\) if

\[
\lim_{x\to c}f(x)=f(c).
\]

\subsubsection{(b) Right-continuity}\label{b-right-continuity}

\(f\) is right-continuous at \(c\) (or continuous from the right) if

\[
\lim_{x\to c^+}f(x)=f(c).
\]

\subsubsection{(c) Left-continuity}\label{c-left-continuity}

\(f\) is left-continuous at \(c\) (or continuous from the left) if

\[
\lim_{x\to c^-}f(x)=f(c).
\]

\section{Additional examples from the
lectures}\label{additional-examples-from-the-lectures}

\subsection{Uniqueness of the limit}\label{uniqueness-of-the-limit}

\textbf{Proposition.} A function cannot have two different limits at the
same point. If \(\lim_{x\to c}f(x)=L\) and \(\lim_{x\to c}f(x)=K\), then
\(L=K\).

\textbf{Proof.} Suppose \(L\ne K\) and choose \(\varepsilon=|L-K|/3\).
The two \(\varepsilon\)-neighborhoods of \(L\) and \(K\) are disjoint,
but the definition of limit would force \(f(x)\) to lie in both for
\(x\) sufficiently close to \(c\). This contradiction proves \(L=K\).

\subsection{\texorpdfstring{A useful \(\varepsilon\)-\(\delta\)
example}{A useful \textbackslash varepsilon-\textbackslash delta example}}\label{a-useful-varepsilon-delta-example}

To prove \(\lim_{x\to2}x^2=4\), observe

\[
|x^2-4|=|x-2||x+2|.
\]

If \(|x-2|<1\), then \(|x+2|<5\). Thus it is enough to take
\(\delta=\min\{1,\varepsilon/5\}\).

\subsection{Power and root rules}\label{power-and-root-rules}

The limit laws imply

\[
\lim_{x\to c}[f(x)]^n=\left(\lim_{x\to c}f(x)\right)^n.
\]

If \(f(x)\ge0\) near \(c\) and the right-hand limit is nonnegative, then

\[
\lim_{x\to c}\sqrt[n]{f(x)}=\sqrt[n]{\lim_{x\to c}f(x)}.
\]

For example,

\[
\lim_{x\to4}\sqrt{x}=2,
\qquad
\lim_{x\to1}\frac{x^3+2x}{x+1}=\frac32.
\]

\subsection{One-sided limit example}\label{one-sided-limit-example}

For \(f(x)=|x|/x\),

\[
\lim_{x\to0^-}f(x)=-1,\qquad \lim_{x\to0^+}f(x)=1,
\]

so the two-sided limit does not exist. This is the simplest model of a
jump discontinuity.

\subsection{Continuity test and types of
discontinuity}\label{continuity-test-and-types-of-discontinuity}

At an interior point \(c\), continuity requires all three conditions:

\begin{enumerate}
\def\labelenumi{\arabic{enumi}.}
\item
  \(f(c)\) is defined;
\item
  \(\lim_{x\to c}f(x)\) exists;
\item
  \(\lim_{x\to c}f(x)=f(c)\).
\end{enumerate}

For the piecewise function in Example 1.1.19, the point \(0\) is a jump
discontinuity because the one-sided limits are \(0\) and \(1\). At
\(1\), the two-sided limit is \(1\) but \(f(1)=2\), so the discontinuity
is removable: changing the value at \(1\) to \(1\) would make the
function continuous there.

\subsection{Riemann-type example}\label{riemann-type-example}

On \([0,1]\), put \(R(p/q)=1/q\) when \(p/q\) is in lowest terms with
\(q>0\), and put \(R(x)=0\) for irrational \(x\). In particular
\(R(0)=R(1)=1\). This is also called Thomae's function. The lecture
extends it periodically outside \([0,1]\).

\textbf{At a rational point.} If \(c=p/q\), irrational points occur
arbitrarily close to \(c\) and have value zero, whereas \(R(c)=1/q\).
Taking \(\varepsilon=1/(2q)\) disproves continuity. At \(0\) or \(1\),
use points from inside the interval.

\textbf{At an irrational point.} Fix an irrational \(c\in(0,1)\) and
\(\varepsilon>0\). Choose an integer \(N>1/\varepsilon\). There are only
finitely many reduced fractions in \([0,1]\) with denominator at most
\(N\); call this finite set \(E_N\). None equals \(c\), so the distance
from \(c\) to \(E_N\) is positive. Choose \(\delta\) smaller than that
distance and than \(c\) and \(1-c\).

If \(|x-c|<\delta\), either \(x\) is irrational and \(R(x)=0\), or its
reduced denominator is greater than \(N\) and \(R(x)<1/N<\varepsilon\).
Since \(R(c)=0\), this is exactly continuity at \(c\).

The useful idea is that the tall spikes form a finite set once a height
is fixed. Density of the rationals by itself does not rule out
continuity.

\subsection{Continuity on an interval and continuous
extension}\label{continuity-on-an-interval-and-continuous-extension}

Continuity on \([a,b]\) means ordinary continuity on \((a,b)\),
right-continuity at \(a\), and left-continuity at \(b\). Use the
corresponding endpoint convention for a half-open interval.

If \(f\) is continuous on an interval except at \(c\) and
\(\lim_{x\to c}f(x)=L\) exists, defining \(f(c)=L\) gives its unique
continuous extension at \(c\). For example, the function \(\sin x/x\)
extends to a continuous function at zero by assigning the value \(1\).
An infinite or oscillating limit cannot be repaired by changing just one
function value.

\subsection{Review problems}\label{review-problems}

\begin{enumerate}
\def\labelenumi{\arabic{enumi}.}
\item
  Prove \(\lim_{x\to2}x^2=4\) without quoting the polynomial limit law.
  Explain why the restriction \(|x-2|<1\) is useful.
\item
  Give functions whose sum has a limit but neither summand does. Then
  give an example showing that strict inequalities need not survive a
  limit.
\item
  Compare \(\sin(1/x)\), \(x\sin(1/x)\) and \(x^2\sin(1/x)\) near zero.
  Which can be extended continuously? Differentiability will be a
  separate question in Week 3.
\end{enumerate}

\subsection{Lecture sources}\label{lecture-sources}

Xing-Bin Pan, MAT1011 lecture slides, Lectures 1--4: limits and limit
language; limit laws and the geometric trigonometric limit; continuity
and the Riemann function (Lecture 4, pp.20--29). The original diagrams
and lecture numbering are retained. Course assessment recorded in the
original draft: class 5\%, homework 15\%, quizzes 20\%, midterm 20\%,
final 40\%.
\end{document}
